% !TEX root = ../main.tex
\chapter{Data Visualization}
\label{ch:visualisation}

\begin{intuition}
Visualization is the data scientist's first tool. A good chart reveals structures, anomalies, and relationships that no table of numbers can show as clearly. As John Tukey said: ``The greatest value of a picture is when it forces us to notice what we never expected to see.''
\end{intuition}

% ============================================================
\section{The Grammar of Graphics}
\index{grammar of graphics}

The \textbf{grammar of graphics}\index{grammar of graphics}, proposed by Leland Wilkinson, decomposes any chart into layers:
\begin{enumerate}
    \item \textbf{Data}: the source dataset.
    \item \textbf{Aesthetics}: the variables mapped to axes, colors, sizes.
    \item \textbf{Geometries}: the type of representation (points, bars, lines).
    \item \textbf{Facets}: the partitioning into sub-plots.
    \item \textbf{Statistics}: the transformations applied (counting, averaging).
    \item \textbf{Coordinates}: the coordinate system (Cartesian, polar).
    \item \textbf{Theme}: the general appearance (fonts, grid, background).
\end{enumerate}

Seaborn and Matplotlib implement these concepts in Python.

% ============================================================
\section{Matplotlib: The Fundamentals}
\index{Matplotlib}

\subsection{Line Plot}

\begin{codeblock}[title=\textbf{Python}]
\begin{minted}{python}
import matplotlib.pyplot as plt
import numpy as np

x = np.linspace(0, 2 * np.pi, 100)
y_sin = np.sin(x)
y_cos = np.cos(x)

fig, ax = plt.subplots(figsize=(7, 4))
ax.plot(x, y_sin, label='sin(x)', color='steelblue', linewidth=2)
ax.plot(x, y_cos, label='cos(x)', color='coral', linestyle='--')
ax.set_xlabel('x')
ax.set_ylabel('y')
ax.set_title('Trigonometric Functions')
ax.legend()
ax.grid(True, alpha=0.3)
plt.tight_layout()
plt.savefig('trigo.pdf', dpi=150)
plt.show()
\end{minted}
\end{codeblock}

\begin{codeoutput}
A plot with two curves (sine in blue, cosine in dashed coral) over $[0, 2\pi]$, with legend, grid, and titles.
\end{codeoutput}

\subsection{Scatter Plot and Histogram}

\begin{codeblock}[title=\textbf{Python}]
\begin{minted}{python}
import seaborn as sns

tips = sns.load_dataset('tips')

fig, axes = plt.subplots(1, 2, figsize=(12, 5))

# Scatter plot
axes[0].scatter(tips['total_bill'], tips['tip'],
                alpha=0.6, c='teal', edgecolors='white')
axes[0].set_xlabel('Total Bill (\$)')
axes[0].set_ylabel('Tip (\$)')
axes[0].set_title('Tip vs Total Bill')

# Histogram
axes[1].hist(tips['total_bill'], bins=20, color='steelblue',
             edgecolor='white')
axes[1].set_xlabel('Total Bill (\$)')
axes[1].set_ylabel('Frequency')
axes[1].set_title('Distribution of Total Bills')

plt.tight_layout()
plt.show()
\end{minted}
\end{codeblock}

\begin{codeoutput}
Two sub-plots: on the left a scatter plot showing a positive relationship between bill and tip; on the right a slightly skewed histogram (right tail) of total bills.
\end{codeoutput}

\subsection{Bar Chart}

\begin{codeblock}[title=\textbf{Python}]
\begin{minted}{python}
moyenne_jour = tips.groupby('day')['tip'].mean().sort_values()

fig, ax = plt.subplots(figsize=(6, 4))
ax.bar(moyenne_jour.index, moyenne_jour.values,
       color=['#4c72b0', '#55a868', '#c44e52', '#8172b2'])
ax.set_xlabel('Day')
ax.set_ylabel('Average Tip (\$)')
ax.set_title('Average Tip by Day')
for i, v in enumerate(moyenne_jour.values):
    ax.text(i, v + 0.05, f'{v:.2f}', ha='center', fontsize=10)
plt.tight_layout()
plt.show()
\end{minted}
\end{codeblock}

\begin{codeoutput}
Colored bar chart with annotated values: Fri~2.73, Thur~2.77, Sat~2.99, Sun~3.26.
\end{codeoutput}

% ============================================================
\section{Seaborn: Statistical Visualization}
\index{Seaborn}

\subsection{Distributions: histplot and violinplot}

\begin{codeblock}[title=\textbf{Python}]
\begin{minted}{python}
fig, axes = plt.subplots(1, 2, figsize=(12, 5))

sns.histplot(data=tips, x='total_bill', hue='sex', kde=True,
             ax=axes[0], palette='Set2')
axes[0].set_title('Distribution by Sex')

sns.violinplot(data=tips, x='day', y='total_bill',
               hue='sex', split=True, ax=axes[1],
               palette='pastel')
axes[1].set_title('Violin Plot by Day and Sex')

plt.tight_layout()
plt.show()
\end{minted}
\end{codeblock}

\begin{codeoutput}
Left: overlapping histograms with KDE curves for male and female. Right: violin plots split by sex for each day, showing that males have slightly higher bills.
\end{codeoutput}

\subsection{Box Plots and Heatmap}

\begin{codeblock}[title=\textbf{Python}]
\begin{minted}{python}
fig, axes = plt.subplots(1, 2, figsize=(13, 5))

sns.boxplot(data=tips, x='day', y='tip', hue='smoker',
            ax=axes[0], palette='coolwarm')
axes[0].set_title('Tips by Day and Smoker Status')

corr = tips[['total_bill', 'tip', 'size']].corr()
sns.heatmap(corr, annot=True, fmt='.2f', cmap='YlOrRd',
            ax=axes[1], vmin=0, vmax=1)
axes[1].set_title('Correlation Matrix')

plt.tight_layout()
plt.show()
\end{minted}
\end{codeblock}

\begin{codeoutput}
Box plots showing similar distributions between smokers and non-smokers. Heatmap: correlation total\_bill/tip = 0.68, total\_bill/size = 0.60, tip/size = 0.49.
\end{codeoutput}

\subsection{Pairplot and relplot}

\begin{codeblock}[title=\textbf{Python}]
\begin{minted}{python}
iris = sns.load_dataset('iris')

g = sns.pairplot(iris, hue='species', palette='husl',
                 diag_kind='kde', height=2.2)
g.figure.suptitle('Pairplot of the Iris Dataset', y=1.02)
plt.show()
\end{minted}
\end{codeblock}

\begin{codeoutput}
A $4 \times 4$ matrix of plots: diagonal = KDE densities by species, off-diagonal = scatter plots. \emph{Setosa} is clearly distinguishable from the other two species.
\end{codeoutput}

\begin{codeblock}[title=\textbf{Python}]
\begin{minted}{python}
sns.relplot(data=tips, x='total_bill', y='tip',
            hue='smoker', style='sex', col='time',
            height=4, aspect=1.2, palette='Set1')
plt.show()
\end{minted}
\end{codeblock}

\begin{codeoutput}
Two panels (Lunch / Dinner) with scatter plots colored by smoker status and differentiated by marker shape according to sex.
\end{codeoutput}

\subsection{catplot}

\begin{codeblock}[title=\textbf{Python}]
\begin{minted}{python}
sns.catplot(data=tips, x='day', y='total_bill', hue='sex',
            kind='box', col='time', palette='muted',
            height=4, aspect=1)
plt.show()
\end{minted}
\end{codeblock}

\begin{codeoutput}
Four groups of box plots organized in two columns (Lunch, Dinner), compared by sex and by day.
\end{codeoutput}

% ============================================================
\section{Choosing the Right Chart Type}
\index{chart selection}

\begin{center}
\begin{tikzpicture}[
    node distance=1.4cm and 2.2cm,
    decision/.style={diamond, draw=teal!80, fill=teal!10, text width=2.8cm, align=center, inner sep=2pt, font=\small},
    block/.style={rectangle, draw=steelblue!80, fill=steelblue!10, rounded corners, text width=2.8cm, align=center, minimum height=0.9cm, font=\small},
    arrow/.style={->, thick, >=stealth, color=gray!70}
]
\node[decision] (start) {Variable type?};
\node[decision, below left=1.5cm and 2.5cm of start] (cat) {Objective?};
\node[decision, below right=1.5cm and 2.5cm of start] (num) {Objective?};

\node[align=center,block, below left=1.3cm and 0.8cm of cat] (bar) {Bar chart\\Counting};
\node[align=center,block, below right=1.3cm and 0.8cm of cat] (catcomp) {Boxplot\\Violin};

\node[align=center,block, below left=1.3cm and 0.8cm of num] (dist) {Histogram\\KDE};
\node[align=center,block, below=1.3cm of num] (rel) {Scatter\\Line};
\node[align=center,block, below right=1.3cm and 0.8cm of num] (multi) {Heatmap\\Pairplot};

\draw[arrow] (start) -- node[above left, font=\scriptsize] {Categorical} (cat);
\draw[arrow] (start) -- node[above right, font=\scriptsize] {Numerical} (num);
\draw[arrow] (cat) -- node[left, font=\scriptsize] {Counting} (bar);
\draw[arrow] (cat) -- node[right, font=\scriptsize] {Comparison} (catcomp);
\draw[arrow] (num) -- node[left, font=\scriptsize] {Distribution} (dist);
\draw[arrow] (num) -- node[right, font=\scriptsize] {Relationship} (rel);
\draw[arrow] (num) -- node[right, font=\scriptsize] {Multivariate} (multi);
\end{tikzpicture}
\end{center}

% ============================================================
\section{Color Palettes and Saving}
\index{color palette}

\begin{codeblock}[title=\textbf{Python}]
\begin{minted}{python}
# Palettes available in Seaborn
palettes = ['deep', 'muted', 'pastel', 'bright',
            'dark', 'colorblind']
sns.set_palette('colorblind')  # recommended for accessibility

# High-resolution saving
fig, ax = plt.subplots()
sns.histplot(tips['tip'], bins=15, ax=ax, color='steelblue')
fig.savefig('distribution_tip.png', dpi=300, bbox_inches='tight')
fig.savefig('distribution_tip.pdf', bbox_inches='tight')
\end{minted}
\end{codeblock}

\begin{bestpractice}
Always use the \py{'colorblind'} palette to ensure the accessibility of your charts. Save in PDF for reports and in PNG (300 dpi) for presentations.
\end{bestpractice}

% ============================================================
\section{Common Visualization Mistakes}
\index{visualization mistakes}

\begin{warning}
\begin{itemize}
    \item \textbf{Truncated axes}: not starting the $y$-axis at zero in a bar chart exaggerates differences.
    \item \textbf{Misleading scales}: using two $y$-axes with different scales can suggest false correlations.
    \item \textbf{3D charts}: perspective distorts the perception of values. Prefer 2D.
    \item \textbf{Too many categories}: a pie chart with 15 slices is unreadable. Use a bar chart instead.
    \item \textbf{Missing legend and titles}: every chart should be understandable on its own.
\end{itemize}
\end{warning}

% ============================================================
\section{Exercises}

\begin{exercise}[$\star$]
Load the \py{tips} dataset and create a histogram of the \py{tip} variable with 20 bins. Add a title, axis labels, and a KDE curve overlay using Seaborn.
\end{exercise}

\begin{exercise}[$\star$]
Create a horizontal bar chart showing the number of meals per day in the \py{tips} dataset. Sort the bars in descending order.
\end{exercise}

\begin{exercise}[$\star\star$]
Using the \py{iris} dataset, create a figure with 4 sub-plots ($2 \times 2$) showing the distribution of each numerical variable, colored by species. Use \py{sns.histplot} with \py{hue='species'}.
\end{exercise}

\begin{exercise}[$\star\star$]
Create a \py{sns.catplot} of type \py{'violin'} showing the distribution of \py{total\_bill} by \py{day}, faceted by \py{time}, with color determined by \py{sex}. Interpret the observed differences.
\end{exercise}

\begin{exercise}[$\star\star\star$]
Reproduce the following chart: a two-panel figure. The left panel shows a scatter plot of \py{total\_bill} vs \py{tip} with a regression line (\py{sns.regplot}). The right panel shows the residuals of this regression (compute them with NumPy: \py{np.polyfit} and \py{np.polyval}). Comment on the quality of the fit.
\end{exercise}

\begin{exercise}[$\star\star\star$]
Take the \py{iris} dataset and create a heatmap of the correlation matrix \emph{per species} (three heatmaps in a $1 \times 3$ figure). Compare the correlation structures across species.
\end{exercise}

% ============================================================
\begin{keyformulas}
\textbf{Chapter Summary -- Data Visualization}
\begin{itemize}
    \item \textbf{Matplotlib}: \py{plt.subplots()}, \py{ax.plot()}, \py{ax.scatter()}, \py{ax.hist()}, \py{ax.bar()}.
    \item \textbf{Seaborn (distributions)}: \py{sns.histplot()}, \py{sns.boxplot()}, \py{sns.violinplot()}.
    \item \textbf{Seaborn (relationships)}: \py{sns.scatterplot()}, \py{sns.relplot()}, \py{sns.heatmap()}.
    \item \textbf{Seaborn (categorical)}: \py{sns.catplot()}, \py{sns.countplot()}, \py{sns.barplot()}.
    \item \textbf{Seaborn (multivariate)}: \py{sns.pairplot()}.
    \item \textbf{Customization}: \py{set\_xlabel()}, \py{set\_title()}, \py{legend()}, \py{grid()}.
    \item \textbf{Saving}: \py{fig.savefig('name.pdf', dpi=300, bbox\_inches='tight')}.
    \item \textbf{Golden rule}: choose the chart suited to the variable type and the analysis objective.
\end{itemize}
\end{keyformulas}
